% Compile from this directory:
% latexmk -pdf -outdir=/tmp/shifted-gaussian shifted-gaussian.tex
\documentclass[11pt]{article}
\usepackage[T1]{fontenc}
\usepackage{lmodern}
\usepackage[a4paper,margin=27mm]{geometry}
\usepackage{microtype}
\usepackage{amsmath,amssymb,amsthm,mathtools}
\usepackage{url}
\usepackage[numbers,sort&compress]{natbib}
\usepackage[hidelinks]{hyperref}
\hypersetup{pdftitle={Centered ellipsoid fitting for non-centered Gaussian data: proof and Rocq formalization}}
\setlength{\emergencystretch}{2em}
\allowdisplaybreaks[1]
\newcommand{\R}{\mathbb R}
\newcommand{\Pbb}{\mathbb P}
\newcommand{\E}{\mathbb E}
\newcommand{\Sym}{\mathbb S}
\DeclareMathOperator{\tr}{tr}
\DeclareMathOperator{\cond}{cond}
\DeclareMathOperator{\Var}{Var}
\newcommand{\opnorm}[1]{\lVert#1\rVert_{\mathrm{op}}}
\newcommand{\norm}[1]{\lVert#1\rVert}
\newcommand{\Fit}{\mathcal F}
\newcommand{\Copt}{\mathcal C}
\newcommand{\OP}{O_{\Pbb}}
\newcommand{\oP}{o_{\Pbb}}
\newtheorem{theorem}{Theorem}[section]
\newtheorem{lemma}[theorem]{Lemma}
\newtheorem{proposition}[theorem]{Proposition}
\theoremstyle{remark}
\newtheorem{remark}[theorem]{Remark}
\begin{document}
\begin{center}
  {\LARGE Centered ellipsoid fitting for non-centered Gaussian data\par}
  \smallskip
  {\large Thresholds, condition numbers, and fitting error\par}
  \medskip
  16 October 2026
\end{center}

\begin{abstract}
For independent $\mathcal N(m_d,I_d)$ points with
$\norm{m_d}=o(\sqrt d)$, the threshold for each prescribed condition
number and the limiting optimal condition number equal their standard
Gaussian values. Without any restriction on the shift, the
unrestricted PSD threshold remains $1/4$. Below it, the optimal
condition number is of order $1+\norm{m_d}/\sqrt d$; above it, one event
excludes every quadric with nonzero PSD quadratic part, whatever its
linear and constant terms. At every positive density the optimal
normalized squared fitting error has the unshifted Gaussian limit.
These centered-fitting statistics also apply exactly to every mixture
of $\mathcal N(m_d,I_d)$ and $\mathcal N(-m_d,I_d)$.

The nonexistence proof corrects a dual certificate using independent
reserve points to annihilate constants and linear functions. It extends
to centrally symmetric admissible mixtures of product laws. For these
mixtures, positive definite fits below the unshifted threshold also
exist for arbitrary shifts supported on $o(d)$ coordinates. The general
mixture inputs and the conclusions of this note are proved in Rocq. The
conclusions are seven public propositions, each proved by a theorem;
their statements are frozen with the repository's checker, and the
formalization is released under the tag \texttt{shifted}.
Appendix~\ref{app:sources} names them and records several differences
between the proofs, including a sparse-shift existence theorem that
does not require symmetry. In particular, two supporting assertions of
the text are not formalized, and the formal proof needs neither: the
almost-sure invertibility of the reserve Gram matrix in
Lemma~\ref{lem:reserve-correction} and the constant of Tropp's matrix
Rademacher bound in~\eqref{eq:reserve-rademacher}. No claim of novelty
is made.
\end{abstract}

\section{Statement and scope}

For a sample $X=(X_1,\ldots,X_n)$ in $\R^d$, let
\[
 \Fit_{n,d}(C)=\{X:\exists Q\succ0,\quad
       X_i^{\mathsf T}QX_i=1\ (1\le i\le n),\quad \cond(Q)\le C\},
\]
where $\cond(Q)=\lambda_{\max}(Q)/\lambda_{\min}(Q)$, and set
\[
 \Copt(X)=\inf\{C\ge1:X\in\Fit_{n,d}(C)\},\qquad \inf\varnothing=+\infty.
\]
The fitting events and $\Copt$ are Borel measurable; the precise formal
sources are listed in Appendix~\ref{app:sources}. The ellipsoids remain
centered at the origin.

Write $\alpha_C=\alpha(C,3)$ for the Gaussian threshold. For completeness,
with
\[
 \nu(dz)=\pi^{-1}\sqrt{2-z^2}\,
       \mathbf1_{[-\sqrt2,\sqrt2]}(z)\,dz,\qquad
 m(t)=\int(z-t)_+\,\nu(dz),\quad v(t)=\int(z-t)_+^2\,\nu(dz),
\]
it is given by
\begin{equation}\label{eq:curve}
 m(-t_C)=C\,m(Ct_C),\qquad
 0<t_C<\sqrt2/C,\qquad
 \alpha_C=\tfrac12-v(-t_C)-v(Ct_C).
\end{equation}
The root is unique, and $C\mapsto\alpha_C$ is continuous and strictly
increasing from $(1,\infty)$ onto $(0,1/4)$. Its inverse is denoted by
$C_\star(\alpha)=C_\star(\alpha,3)$. Appendix~\ref{app:sources} gives
the formal sources for the root, the curve properties and the inverse.

\begin{theorem}[Non-centered Gaussian data]\label{thm:shifted}
Let $m_d\in\R^d$ be deterministic and put $a_d=\norm{m_d}$.
Suppose
\begin{equation}\label{eq:growth}
 \frac{a_d}{\sqrt d}\longrightarrow0,
 \qquad \frac{n(d)}{d^2}\longrightarrow\alpha>0.
\end{equation}
Let $X_i=Z_i+m_d$, where the $Z_i$ are independent
$\mathcal N(0,I_d)$ vectors. For every fixed $C>1$,
\begin{equation}\label{eq:threshold}
 \Pbb\{X\in\Fit_{n(d),d}(C)\}\longrightarrow
 \begin{cases}
  1,&\alpha<\alpha_C,\\
  0,&\alpha>\alpha_C.
 \end{cases}
\end{equation}
If $0<\alpha<1/4$, then
\begin{equation}\label{eq:optimum}
 \Copt(X)\xrightarrow{\Pbb}C_\star(\alpha).
\end{equation}
\end{theorem}

The point law has mean $m_d$ and covariance $I_d$, so it need not
satisfy the centering assumption of the existing mixture theorem.
Since $\alpha_C\uparrow1/4$, the result gives positive definite fits
at every $0<\alpha<1/4$ by taking $C$ large enough. Theorem~\ref{thm:unrestricted}
also gives bounded condition numbers for shifts of order $\sqrt d$ and
matching growth bounds for larger shifts. It does not identify the
precise condition-number curve when $a_d/\sqrt d$ has a nonzero limit.
No assertion is made at $\alpha=\alpha_C$.

\begin{remark}[The sign mixture gives the same fitting problem]
\label{rem:sign-invariance}
For any $w_d\in[0,1]$, the conclusions of Theorem~\ref{thm:shifted} also
hold for independent points with law
\[
 w_d\mathcal N(m_d,I_d)+(1-w_d)\mathcal N(-m_d,I_d).
\]
In fact, the equality of fitting statistics is exact at every $n,d$.
Write these points as $\widetilde X_i=Z_i+\sigma_i m_d$, with independent
signs satisfying $\Pbb(\sigma_i=1)=w_d$, independent of the $Z_i$.
Then $Y_i=\sigma_i\widetilde X_i=\sigma_i Z_i+m_d$ are independent
$\mathcal N(m_d,I_d)$ points, while
$Y_iY_i^{\mathsf T}=\widetilde X_i\widetilde X_i^{\mathsf T}$.
Thus the entire sets of fitting matrices coincide samplewise, as do
their optimal condition numbers and normalized fitting errors. This
includes $w_d=0,1$ and does not
assert equality of the point distributions. The identity applies to
positive semidefinite as well as positive definite fitting matrices and
requires no bound on $\norm{m_d}$.
\end{remark}

\begin{theorem}[The unrestricted threshold for arbitrary shifts]
\label{thm:unrestricted}
Let $n(d)/d^2\to\alpha\in(0,\infty)$, let $m_d\in\R^d$ be arbitrary
and deterministic, and let $X_i=Z_i+m_d$ for independent
$Z_i\sim\mathcal N(0,I_d)$. There is no restriction on $\norm{m_d}$.
\begin{enumerate}
\item If $\alpha<1/4$, then with probability tending to one there is
      $Q\succ0$ such that $X_i^{\mathsf T}QX_i=1$ for every $i$.
      More precisely, there are constants $0<c_\alpha\le C_\alpha<\infty$
      such that, with probability tending to one,
      \begin{equation}\label{eq:condition-growth}
       c_\alpha\left(1+\frac{\norm{m_d}}{\sqrt d}\right)
       \le\Copt(X)\le
       C_\alpha\left(1+\frac{\norm{m_d}}{\sqrt d}\right).
      \end{equation}
\item If $\alpha>1/4$, then with probability tending to one there is
      no $Q\succeq0$ such that $X_i^{\mathsf T}QX_i=1$ for every $i$.
      More strongly, there are Borel sets $\Omega_d$ of Gaussian samples
      with $\Pbb(\Omega_d)\to1$ such that, on $\Omega_d$, no
      $Q\succeq0$, $Q\ne0$, $\ell\in\R^d$, and $c\in\R$ satisfy
      \begin{equation}\label{eq:no-quadric}
       Z_i^{\mathsf T}QZ_i+2\ell^{\mathsf T}Z_i+c=0
       \qquad(1\le i\le n).
      \end{equation}
      In particular, no PSD fit exists for $Z_i+m$, simultaneously
      for every $m\in\R^d$. No ellipsoid, whatever its center, passes
      through all the points.
\end{enumerate}
The existence and nonexistence conclusions for centered fits, and the
condition-number bounds, also hold for independent points with law
$w_d\mathcal N(m_d,I_d)+(1-w_d)\mathcal N(-m_d,I_d)$, for every
$w_d\in[0,1]$.
\end{theorem}

Section~\ref{sec:unrestricted} proves this theorem and identifies the
limiting fitting error in Theorem~\ref{thm:shifted-error}.
Section~\ref{sec:symmetric-mixtures} extends the unrestricted
nonexistence argument to symmetric non-Gaussian mixtures.
Exact fitting at the critical density, general component means and
covariances, and existence for dense shifts of non-Gaussian data are
not settled here.

\section{Inputs and deterministic perturbation}

Here $\Sym_p$ denotes the real symmetric $p\times p$ matrices,
$\norm\cdot$ is the Euclidean norm, and logarithms are natural.
The general mixture condition-number and interpolation theorems are
proved in Rocq. We use their elementary Gaussian specializations.
For a standard normal scalar $G_0$,
\begin{equation}\label{eq:normal-admissibility}
 \E G_0=0,\qquad \E G_0^2=1,\qquad \E G_0^4=3,
 \qquad \E e^{G_0^2/4}=\sqrt2\le2.
\end{equation}
Hence the one-component Gaussian family is admissible with $K=2$ and
fourth moment three. The Rocq development proves these scalar facts and
exports the Gaussian specializations used below, in
\path{GaussianInputs.v} and \path{GaussianExports.v} of
\path{extension/Shifted/}; see Appendix~\ref{app:sources}.

\paragraph{Gaussian fitting.}
If $g_1,\ldots,g_N$ are independent $\mathcal N(0,I_p)$ and
$N/p^2\to\beta>0$, then for fixed $D>1$ the probability of
$\Fit_{N,p}(D)$ tends to one for $\beta<\alpha_D$, and to zero for
$\beta>\alpha_D$. This follows by applying
\path{condition_statement_holds} to the family just described.

\paragraph{Simultaneous interpolation.}
If $0<\beta<1/2$ and $N/p^2\to\beta$, there is $L_\beta>0$ such that,
outside Borel exceptional sets of probability tending to zero, the
Gaussian sample has the following property: for every $y\in\R^N$ there
exist $\Delta\in\Sym_p$ and $b\in\R$ with
\begin{align}
 g_i^{\mathsf T}\Delta g_i-\tr\Delta
       &=\sqrt p\,(y_i+b)\qquad(1\le i\le N),\label{eq:interpolation}\\
 \opnorm{\Delta}
       &\le L_\beta\left(\frac{\norm y_2}{p}
          +\frac{(\log p)^4}{\sqrt p}\norm y_\infty\right).
          \label{eq:interpolation-bound}
\end{align}
This is the Gaussian specialization of
\path{mixture_interpolation_holds} in
\path{extension/Mixture/Interpolation.v}, obtained from
\path{mixture_small_residual}; its source is the interpolation theorem
of Koehler--Sohn~\cite[Theorem 1.9]{KoehlerSohn2026}.
The formal statement gives an exceptional set of probability at most
each prescribed $\epsilon>0$ in every sufficiently large dimension.
A diagonal choice gives the sets used here.

The conclusion holds simultaneously for every $y$. We use it on the
same sample as the Gaussian fitting event. All probability bounds use
the Borel exceptional sets, without asserting measurability of the full
interpolation-property event or choosing a measurable correction matrix.

\begin{lemma}[Stability under a small correction]\label{lem:stability}
Suppose $A\succ0$, $\tr A=p$, and $\cond(A)\le D$. Then
$\lambda_{\min}(A)\ge1/D$ and $\opnorm A\le D$. If $E$ is symmetric
and $\opnorm E\le\delta<1/D$, then $A+E\succ0$ and
\begin{equation}\label{eq:condition-stability}
 \cond(A+E)\le D+\frac{(D+1)\delta}{1/D-\delta}.
\end{equation}
Also, $\operatorname{diag}(1,A)$ has trace $p+1$ and condition number
at most $D$.
\end{lemma}
\begin{proof}
The mean eigenvalue of $A$ is one, so
$\lambda_{\min}(A)\le1\le\lambda_{\max}(A)$ and the assertions about
$A$ and $\operatorname{diag}(1,A)$ follow. Writing
$\ell=\lambda_{\min}(A)$, the extremal eigenvalue bounds give
\[
 \cond(A+E)\le\frac{D\ell+\delta}{\ell-\delta}
   =D+\frac{(D+1)\delta}{\ell-\delta}
   \le D+\frac{(D+1)\delta}{1/D-\delta}.\qedhere
\]
\end{proof}

We will repeatedly rescale a fitting matrix to have trace equal to the
dimension. Its quadratic values then equal some common level $q>0$,
instead of one. Conversely, a positive definite matrix with a common
quadratic value on nonzero observations has a positive common level;
division by that level gives an exact fit without changing its condition
number. This observation removes any need to estimate the scalar $b$
in~\eqref{eq:interpolation}.

\section{Existence below the threshold}
\label{sec:conditioned-existence}

Assume $\alpha<\alpha_C$ and choose $1<C_-<C$ with
$\alpha<\alpha_{C_-}$. Rotate $m_d$ to $a_de_1$ and write
\[
 X_i=(u_i,g_i),\qquad u_i=\xi_i+a_d,
 \qquad p=d-1,
\]
where the $\xi_i$ are standard Gaussian, the $g_i$ are
$\mathcal N(0,I_p)$, and all these variables are independent. This
uses Gaussian rotation invariance and splitting off the first
coordinate. Since $n(d)/p^2\to\alpha$, Gaussian fitting
gives, with probability tending to one, a matrix $A\succ0$ such that
\[
 \tr A=p,\qquad \cond(A)\le C_-,\qquad
 g_i^{\mathsf T}Ag_i=q\quad\hbox{for every }i.
\]
Interpolate on this same Gaussian sample the centered targets
\begin{equation}\label{eq:sat-targets}
 y_i=-\frac{u_i^2-(1+a_d^2)}{\sqrt p}.
\end{equation}
Here
\[
 u_i^2-(1+a_d^2)=\xi_i^2-1+2a_d\xi_i,
 \qquad \Var(u_i^2)=2+4a_d^2.
\]
Since $n(d)=O(d^2)$, the second-moment bound and the Gaussian maximum
bound give
\begin{align}
 \frac{\norm y_2}{p}
    &=\OP\left(\frac{1+a_d}{\sqrt d}\right),\label{eq:sat-l2}\\
 \norm y_\infty
    &=\OP\left(\frac{a_d\sqrt{\log d}+\log d}{\sqrt d}\right).
       \label{eq:sat-linf}
\end{align}
For example, $\E\norm y_2^2=n(d)(2+4a_d^2)/p$ proves the first
estimate by Markov's inequality, and
$\max_i|\xi_i|=\OP(\sqrt{\log d})$ proves the second.
Because $\alpha<1/4<1/2$, interpolation applies. Outside its Borel
exceptional set, there exists a correction with
\[
 \opnorm\Delta\le S_d:=L_\alpha\left(
   \frac{\norm y_2}{p}+\frac{(\log p)^4}{\sqrt p}\norm y_\infty\right).
\]
This upper bound is a measurable random variable, and
\begin{equation}\label{eq:sat-correction}
 S_d=\OP\left(
  \frac{1+a_d}{\sqrt d}
  +\frac{a_d(\log d)^{9/2}}{d}
  +\frac{(\log d)^5}{d}\right)=\oP(1)
\end{equation}
whenever $a_d=o(\sqrt d)$.

Set $Q=\operatorname{diag}(1,A+\Delta)$ and choose
\[
 \delta=\frac{C-C_-}{2C_-(C+1)}<\frac1{C_-}.
\]
Then $C_-+(C_-+1)\delta/(1/C_--\delta)\le C$.
On $\{S_d\le\delta\}$, Lemma~\ref{lem:stability}, applied to
$\operatorname{diag}(1,A)$, gives $Q\succ0$ and $\cond(Q)\le C$.
Moreover,
\begin{align*}
 X_i^{\mathsf T}QX_i
 &=u_i^2+q+\tr\Delta+\sqrt p\,(y_i+b)\\
 &=q+1+a_d^2+\tr\Delta+\sqrt p\,b.
\end{align*}
This is independent of $i$. It is positive because $Q\succ0$ and the
observations are almost surely nonzero. Rescaling $Q$ gives an exact
fit of condition number at most $C$. The Gaussian fitting event, the
complement of the interpolation exceptional set, and $\{S_d\le\delta\}$
have intersection of probability tending to one. This proves the first half
of~\eqref{eq:threshold}.

\section{A uniform residual gap for Gaussian points}
\label{sec:data-gap}

The missing quantitative statement is contained in the proof of
\path{condition_unsat_gap}, rather than in its exported conclusion.
We extract it here, retaining the squared loss that the exact-fit proof
sets to zero. This step concerns the actual quadratic features of
Gaussian points, not just auxiliary Gaussian matrix measurements.

Write
\[
 \mathcal K_d(C)=\{R\in\Sym_d:R\succ0,\ \tr R=d,\ \cond(R)\le C\},
 \qquad W_i=\frac{Z_iZ_i^{\mathsf T}-I_d}{\sqrt d}.
\]
For a tuple $U=(U_1,\ldots,U_n)$ of symmetric matrices define
\[
 \mathcal E_U(R,b)=\frac1n\sum_{i=1}^n
                    (b-\langle U_i,R\rangle)^2,
 \qquad \langle A,B\rangle=\tr(AB).
\]
Let $P_n$ be orthogonal projection onto the vectors in $\R^n$ with
coordinate sum zero. For any point sample $Z$, set
\begin{equation}\label{eq:centered-measurements}
 \mathcal L_Z(R)=\frac1{\sqrt{nd}}\,
                  P_n\bigl((Z_i^{\mathsf T}RZ_i)_{i=1}^n\bigr).
\end{equation}
Subtracting $\tr R$ from every quadratic value has no effect under
$P_n$, so
\begin{equation}\label{eq:loss-identity}
 \norm{\mathcal L_Z(R)}_2^2=\min_{b\in\R}\mathcal E_W(R,b).
\end{equation}

\begin{lemma}[Extracted data gap]\label{lem:data-gap}
Let $C>1$, $\alpha>\alpha_C$, and $n(d)/d^2\to\alpha$.
Choose $C_p>C$ with $\alpha_{C_p}<\alpha$, and put
\begin{equation}\label{eq:gap-constant}
 g=\frac{(\sqrt\alpha-\sqrt{\alpha_{C_p}})^2}{64\alpha}>0.
\end{equation}
For independent $Z_i\sim\mathcal N(0,I_d)$ there are Borel sets
$\Omega_d$ with $\Pbb(\Omega_d)\to1$ such that, on $\Omega_d$,
\begin{equation}\label{eq:data-gap}
 \mathcal E_W(R,b)\ge g/4
 \qquad\text{for every }R\in\mathcal K_d(C),\ b\in\R.
\end{equation}
In particular, $\inf_{R\in\mathcal K_d(C)}
\norm{\mathcal L_Z(R)}_2^2\ge g/4$ there.
\end{lemma}

\begin{proof}
We spell out the estimates from
\path{extension/Condition/Unsat.v} that make the assertion uniform in
$R$ and retain a positive residual. The Gaussian specialization is
justified by~\eqref{eq:normal-admissibility}.

\emph{Cover and regularized objectives.}
The finite cover \path{cond_cover} supplies numbers
$0<\ell_j<1<h_j<C_p\ell_j$, $1\le j\le J$, and a barrier bound $B$,
independent of $d$. Enlarge $B$ to $\max(B,0)$. Every
$R\in\mathcal K_d(C)$ lies in some slice
\[
 \mathcal D_{j,d}=\{S\in\Sym_d:\tr S=d,\ \ell_j I_d\prec S\prec h_jI_d\}
\]
with $\mathcal B_{j,d}(R)\le dB$. Here the normalized barrier is
\[
 \mathcal B_{j,d}(S)=-\log\det(S-\ell_jI_d)
    -\log\det(h_jI_d-S)+d\log((1-\ell_j)(h_j-1)).
\]
It is nonnegative on the slice. Let $\lambda_d$ be the operator norm
of the covariance of $W_i$, so that $0<d\lambda_d\le2$, and put
$\gamma_d=(d\sqrt d)^{-1}$. For positive weights $\beta_j$ chosen below,
the slice value used in the formal proof is
\begin{equation}\label{eq:slice-value}
 \Phi_j(U)=\inf_{\substack{S\in\mathcal D_{j,d}\\b\in\R}}
 \left\{\frac12\mathcal E_U(S,b)
   +\beta_j(\lambda_d\norm S_F^2+b^2)
   +\gamma_d\mathcal B_{j,d}(S)\right\}.
\end{equation}
The norm $\norm\cdot_F$ is Frobenius norm. This is \path{af_value} at
fourth moment three.

\emph{The simultaneous lower bounds.}
For auxiliary independent centered Gaussian matrix measurements $H_i$ with the
same covariance as $W_i$, \path{gaussian_gap} gives
$\mathcal E_H(S,b)\ge g$ throughout $\mathcal K_d(C_p)$ outside
an event of probability tending to zero. Each $\mathcal D_{j,d}$ is
contained in that cone slice. Nonnegativity of the ridge and barrier
therefore gives $\Phi_j(H)\ge g/2$ on this event, and
$\E\Phi_j(H)\ge g/4$ for all large $d$.
The comparison and concentration estimates used in
\path{condition_unsat_gap} then yield, simultaneously for all $j$,
\begin{equation}\label{eq:slice-lower}
 \Phi_j(W)\ge g/4-\beta_j/4
\end{equation}
outside Borel exceptional sets of probability tending to zero.
These are precisely \path{af_center_ge}, \path{af_conc},
\path{sl_err_small}, and \path{sl_prob_small}, applied to the finitely
many fixed slices. None of these events assumes an exact fit.

\emph{The minimizing intercept and the upper bound.}
On the Borel a priori event of \path{mixture_adjoint_bound}, for a
constant $C_a$ and every $z\in\R^n$,
\[
 \left\|\sum_i z_iW_i\right\|_F\le C_a\sqrt d\,\norm z_2.
\]
For each $R\in\mathcal K_d(C)$ choose a covering slice $j$ and set
\[
 b_R=\frac1n\sum_i\langle W_i,R\rangle
       =\langle\overline W,R\rangle,\qquad
 \overline W=\frac1n\sum_i W_i.
\]
The a priori bound and $\norm R_F^2\le h_jd$ give, eventually,
\begin{equation}\label{eq:intercept-bound}
 |b_R|\le C_a\sqrt{d/n}\,\norm R_F,
 \qquad b_R^2\le c_bh_j,\qquad c_b=2C_a^2/\alpha,
\end{equation}
using $n\ge\alpha d^2/2$. Unlike \path{common_value_le}, this argument
does not assume equal feature values: $b_R$ is their empirical mean by
definition.

Use the same parameter choices as the formal proof:
\[
 \theta=\frac{g}{16(1+c_b)},\qquad
 \beta_j=\frac{\theta}{2(2h_j+1)}.
\]
Then
\begin{equation}\label{eq:ridge-bounds}
 \beta_j\le g/32,\qquad
 \beta_j(2h_j+b_R^2)\le g/32,\qquad
 \gamma_d\mathcal B_{j,d}(R)\le B/\sqrt d\le g/32
\end{equation}
for all sufficiently large $d$. Evaluating~\eqref{eq:slice-value} at
$(R,b_R)$, and using $d\lambda_d\le2$, gives
\begin{equation}\label{eq:slice-upper}
 \Phi_j(W)\le\tfrac12\mathcal E_W(R,b_R)+g/16.
\end{equation}
This is exactly the proof of \path{af_value_common} with its loss term
retained instead of set to zero.

Combining~\eqref{eq:slice-lower}--\eqref{eq:slice-upper},
\[
 \tfrac12\mathcal E_W(R,b_R)+g/16
       \ge g/4-g/128=31g/128,
 \qquad \mathcal E_W(R,b_R)\ge23g/64>g/4.
\]
Finally,
$\mathcal E_W(R,b)=\mathcal E_W(R,b_R)+(b-b_R)^2$,
so the bound holds for every $b$. The event $\Omega_d$ is the
intersection of the a priori event and the complements of the finitely
many slice exceptional sets, with the same diagonal choice as in
Section~2. It is Borel and has probability tending to one, independently
of $R$ and $b$. No restriction $\alpha<1/2$ was used.
\end{proof}

\section{Nonexistence and the optimal condition number}

Assume $\alpha>\alpha_C$ and use the full sample of $n=n(d)$ points.
Let $G$ have rows $Z_i^{\mathsf T}$. The standard Gaussian matrix norm
bound gives $\opnorm G=\OP(\sqrt n+\sqrt d)=\OP(d)$.
Expanding $X_i=Z_i+m_d$ and applying $P_n$ removes the common term
$m_d^{\mathsf T}Rm_d$. Uniformly over $R\in\mathcal K_d(C)$,
\begin{equation}\label{eq:residual-comparison}
 \norm{\mathcal L_X(R)-\mathcal L_Z(R)}_2
 =\frac{2\norm{P_nGRm_d}_2}{\sqrt{nd}}
 \le T_d:=\frac{2Ca_d\opnorm G}{\sqrt{nd}}
 =\OP\left(\frac{a_d}{\sqrt d}\right)=\oP(1).
\end{equation}
Here $\opnorm R\le C$ by Lemma~\ref{lem:stability}, and $T_d$ is a
measurable bound independent of the choice of $R$.

A hypothetical fit of $X$ with condition number at most $C$, rescaled
to trace $d$, belongs to $\mathcal K_d(C)$ and satisfies
$\mathcal L_X(R)=0$. On $\Omega_d$ from Lemma~\ref{lem:data-gap},
it would therefore imply $\sqrt g/2\le T_d$. Consequently,
\[
 \Pbb\{X\in\Fit_{n,d}(C)\}
 \le\Pbb(\Omega_d^c)+\Pbb\{T_d\ge\sqrt g/2\}\longrightarrow0.
\]
This proves nonexistence without subsampling or reverse interpolation.
Uniformity in~\eqref{eq:residual-comparison} handles a matrix selected
from the data; no measurable choice of a fitting matrix is needed.

For $0<\alpha<1/4$, the optimum now follows from the threshold and
monotonicity of $\Fit_{n,d}(C)$ in $C$. Given $\eta>0$, existence at
$C_\star(\alpha)+\eta/2$ gives
$\Pbb\{\Copt(X)>C_\star(\alpha)+\eta\}\to0$.
If $C_\star(\alpha)-\eta\le1$, the lower deviation is impossible.
Otherwise choose $D\in(C_\star(\alpha)-\eta,C_\star(\alpha))$.
The definition of an infimum and monotonicity give
\[
 \{\Copt(X)<C_\star(\alpha)-\eta\}
       \subseteq\{X\in\Fit_{n,d}(D)\}.
\]
The right-hand probability tends to zero because $\alpha>\alpha_D$.
This proves~\eqref{eq:optimum} without assuming attainment and completes
Theorem~\ref{thm:shifted}.

\section{Unrestricted fits for arbitrary Gaussian shifts}
\label{sec:unrestricted}

The bound in~\eqref{eq:residual-comparison} uses a fixed condition
number. We now prove Theorem~\ref{thm:unrestricted} by a different
argument. Above $1/4$ we construct a positive definite dual certificate
whose weights annihilate both constants and linear functions. Such a
certificate is unchanged by translation of the observations.

\subsection{Positive definite fits below \texorpdfstring{$1/4$}{1/4}}

Suppose $0<\alpha<1/4$ and choose a fixed $D>1$ with
$\alpha<\alpha_D$. As in Section~\ref{sec:conditioned-existence}, rotate
the shift to $a_de_1$ and write $X_i=(\xi_i+a_d,g_i)$, with $p=d-1$.
With probability tending to one, the Gaussian sample $(g_i)$ admits
$A\succ0$ with $\tr A=p$, $\cond(A)\le D$, and
$g_i^{\mathsf T}Ag_i=q$ independent of $i$.

Fix a sufficiently small $\varepsilon>0$, depending only on $\alpha$
and $D$, and give the first coordinate the positive weight
\[
 s_d=\min\left(1,\frac{\varepsilon\sqrt d}{a_d}\right),
 \qquad s_d=1\quad\text{if }a_d=0.
\]
Interpolate the targets
\[
 y_i=-\frac{s_d\big((\xi_i+a_d)^2-(1+a_d^2)\big)}{\sqrt p}.
\]
Since $s_d\le1$ and $s_da_d\le\varepsilon\sqrt d$, the law of large
numbers for $\xi_i^2$ and $(\xi_i^2-1)^2$ gives
\begin{align}
 \frac{\norm y_2}{p}
 &\le\frac{s_d}{p^{3/2}}\norm{(\xi_i^2-1)_i}_2
       +\frac{2s_da_d}{p^{3/2}}\norm{(\xi_i)_i}_2
 \le2\sqrt\alpha\,\varepsilon+\oP(1),
       \label{eq:weighted-target-l2}\\
 \frac{(\log p)^4}{\sqrt p}\norm y_\infty
 &=\OP\left(\frac{(\log d)^5}{d}
       +\frac{\varepsilon(\log d)^{9/2}}{\sqrt d}\right)=\oP(1).
       \label{eq:weighted-target-infinity}
\end{align}
The second estimate uses $\max_i|\xi_i|=\OP(\sqrt{\log d})$.
In particular, the first bound has a deterministic leading constant;
a tightness estimate alone would not suffice for this choice of $s_d$.
Simultaneous interpolation therefore supplies a symmetric correction
$\Delta$ and an intercept $b_0$ with the measurable upper bound
\[
 \opnorm\Delta\le S'_d:=L_\alpha\left(
 \frac{\norm y_2}{p}+\frac{(\log p)^4}{\sqrt p}\norm y_\infty\right)
 \le2L_\alpha\sqrt\alpha\,\varepsilon+\oP(1).
\]
Choose $0<\varepsilon\le1$ so that
$2L_\alpha\sqrt\alpha\,\varepsilon<1/(4D)$. With probability tending
to one, $\opnorm\Delta\le1/(2D)$, and hence
\[
 (2D)^{-1}I_p\preceq A+\Delta\preceq2D I_p.
\]
The matrix
$Q=\operatorname{diag}(s_d,A+\Delta)$ is then positive definite and
\[
 X_i^{\mathsf T}QX_i
 =q+s_d(1+a_d^2)+\tr\Delta+\sqrt p\,b_0
\]
is independent of $i$. Its common value is positive, so rescaling gives
an exact positive definite fit, with
\begin{equation}\label{eq:weighted-condition-upper}
 \cond(Q)\le\frac{4D^2}{s_d}
 \le\frac{4D^2}{\varepsilon}\left(1+\frac{a_d}{\sqrt d}\right).
\end{equation}
The Borel exceptional-set argument is the same as in
Section~\ref{sec:conditioned-existence}. This proves existence and the
upper bound in~\eqref{eq:condition-growth}, without a growth restriction
on $a_d$.

\paragraph{A matching lower bound.}
For any fixed density $\alpha>0$, let $G$ have rows $Z_i^{\mathsf T}$.
The adjoint bound used in~\eqref{eq:reserve-adjoint}, by duality, gives
on events of probability tending to one a fixed $L>0$, depending only
on $\alpha$, such that
\[
 \norm{(Z_i^{\mathsf T}QZ_i-\tr Q)_i}_2
       \le Ld\,\norm Q_F\qquad(Q\in\Sym_d).
\]
Also $\sigma_{\min}(P_nG)/\sqrt n\to1$ in probability: an orthogonal
change of the row coordinates identifies its nonzero rows with an
$(n-1)\times d$ standard Gaussian matrix. If $Q\succ0$ fits $Z+m$,
centering its fit equations gives
\[
 P_n(Z_i^{\mathsf T}QZ_i)_i=-2P_nGQm.
\]
Consequently, on the same events, simultaneously for every such $m,Q$,
\begin{align*}
 2\sigma_{\min}(P_nG)\lambda_{\min}(Q)\norm m
 &\le2\norm{P_nGQm}_2\\
 &\le Ld\norm Q_F
 \le Ld^{3/2}\lambda_{\max}(Q).
\end{align*}
Canceling $\lambda_{\min}(Q)>0$ gives the scale-invariant bound
\begin{equation}\label{eq:condition-lower}
 \cond(Q)\ge
 \frac{2\sigma_{\min}(P_nG)}{Ld}\,
       \frac{\norm m}{\sqrt d}
 =\left(\frac{2\sqrt\alpha}{L}+\oP(1)\right)
       \frac{\norm m}{\sqrt d}.
\end{equation}
The random coefficient is independent of $m$ and $Q$. Together with
$\cond(Q)\ge1$, this proves the lower bound
in~\eqref{eq:condition-growth} after decreasing a deterministic positive
constant $c_\alpha$. Taking the infimum over fitting matrices requires
no attainment assumption. In particular, for $\alpha<1/4$, shifts of
order $\sqrt d$ admit bounded condition numbers, whereas
$a_d/\sqrt d\to\infty$ forces the optimal condition number to diverge
linearly in $a_d/\sqrt d$, up to fixed multiplicative constants.

\subsection{A Gaussian certificate with a positive margin}

For a sample $Z=(Z_1,\ldots,Z_N)$ define its unrestricted normalized
squared residual by
\begin{equation}\label{eq:psd-residual}
 \Gamma_{N,d}(Z)=
 \min_{\substack{R\succeq0\\\tr R=d}}
       \norm{\mathcal L_Z(R)}_2^2.
\end{equation}
Here $\mathcal L_Z$ is~\eqref{eq:centered-measurements}, with sample
size $N$. For a PSD matrix the quadratic values are nonnegative, so
subtracting their empirical mean is also the minimization over
nonnegative common levels used in the Gaussian residual theorem.
For $N/d^2\to\beta>0$ and standard Gaussian points, the squared-error
theorem of \citet[Theorems 1.2 and 1.7]{KoehlerSohn2026} gives
\begin{equation}\label{eq:gaussian-error-limit}
 \Gamma_{N,d}(Z)\xrightarrow{\Pbb}e_\star(\beta,3),\qquad
 e_\star(\beta,3)=\frac{2}{\beta}(-\omega_\beta)_+^2,
 \qquad v(\omega_\beta)=\beta.
\end{equation}
Here $v$ is defined in Section~1. It is continuous and strictly
decreasing from $+\infty$ to $0$ on $(-\infty,\sqrt2)$, so
$\omega_\beta$ is unique and continuous in $\beta>0$.
Thus $e_\star(\beta,3)$ is continuous, equals zero for
$\beta\le1/4$, and is positive for $\beta>1/4$, since $v(0)=1/4$.
In particular, for $\beta>1/4$ there is $\eta_\beta>0$ such that
\begin{equation}\label{eq:psd-gaussian-gap}
 \Pbb\{\Gamma_{N,d}(Z)\ge\eta_\beta\}\longrightarrow1.
\end{equation}
Unlike Lemma~\ref{lem:data-gap}, this input ranges over the entire PSD
trace slice. The following consequence is the dual identity of
\citet[Proposition A.1]{KoehlerSohn2026}, in our normalization. We give
the construction to make its measurability and sign invariance explicit.

\begin{lemma}[A certificate determined by the outer products]
\label{lem:psd-certificate}
Let $N/d^2\to\beta>1/4$. There is a Borel choice of
$y=y(Z)\in\R^N$, depending only on $(Z_iZ_i^{\mathsf T})_{i\le N}$,
such that $\sum_i y_i=0$, $\norm y_2\le1$, and, for any fixed
$0<c_\beta<\sqrt{\beta e_\star(\beta,3)}$, with probability tending to one,
\begin{equation}\label{eq:certificate-margin}
 M_0:=\sum_{i=1}^N y_iZ_iZ_i^{\mathsf T}
       \succeq c_\beta\sqrt d\,I_d.
\end{equation}
Its linear moment $b=\sum_i y_iZ_i$ satisfies
$\norm b=\OP(\sqrt d)$.
\end{lemma}

\begin{proof}
The image
\[
 \mathcal V_Z=
 \{\mathcal L_Z(R):R\succeq0,\ \tr R=d\}\subset\mathbf1^\perp
\]
is compact and convex. Let $z_*(Z)$ be its unique point of smallest
Euclidean norm. Set $y=z_*/\norm{z_*}$ when $z_*\ne0$ and $y=0$
otherwise. The vector $z_*$ depends continuously on the sample: along a
convergent sequence of samples, compactness of the fixed PSD trace slice
gives subsequential limits of minimizing matrices, and every resulting
image limit is the unique minimizing image at the limiting sample.
Consequently $y$ is Borel. Its definition uses only the outer products.

The projection inequality gives, for every $R\succeq0$ of trace $d$,
\[
 \langle y,\mathcal L_Z(R)\rangle
 \ge\norm{z_*}=\sqrt{\Gamma_{N,d}(Z)}.
\]
This also holds when $z_*=0$. Since $y\perp\mathbf1$, it implies
\[
 \tr(M_0R)\ge\sqrt{Nd}\,\sqrt{\Gamma_{N,d}(Z)}.
\]
Minimizing the left side over the PSD trace slice gives
\begin{equation}\label{eq:certificate-duality}
 \lambda_{\min}(M_0)
 \ge\sqrt{N/d}\,\sqrt{\Gamma_{N,d}(Z)}.
\end{equation}
Equation~\eqref{eq:gaussian-error-limit} proves~\eqref{eq:certificate-margin}
with every constant $c_\beta$ in the stated range.

To bound $b$, represent $Z_i=\varepsilon_iU_i$ with independent
standard Gaussian $U_i$ and independent Rademacher signs
$\varepsilon_i$. The weights $y_i$ depend on the $U_iU_i^{\mathsf T}$
and hence not on these signs. Conditional on $(U_i)_i$,
\[
 \E_\varepsilon\left\|\sum_i y_i\varepsilon_iU_i\right\|^2
 =\sum_i y_i^2\norm{U_i}^2
 \le\max_i\norm{U_i}^2.
\]
The maximum is at most $Cd$ with probability tending to one because
$N=O(d^2)$. Conditional Markov's inequality now gives
$\norm b=\OP(\sqrt d)$.
\end{proof}

\subsection{Correcting the linear moment on an independent sample}

\begin{lemma}[Correction with a small quadratic contribution]
\label{lem:reserve-correction}
Suppose $k=k(d)$ satisfies $k/d^2\to\gamma>0$. Let $H$ be a $k\times d$
matrix with independent standard Gaussian rows $H_j^{\mathsf T}$, and
let $b\in\R^d$ be independent of $H$, with $\norm b=\OP(\sqrt d)$.
Put
\[
 P_k=I_k-\frac1k\mathbf1\mathbf1^{\mathsf T},\qquad
 K=H^{\mathsf T}P_kH,\qquad v=K^{-1}b,\qquad r=P_kHv.
\]
For all sufficiently large $d$, $K$ is invertible almost surely. This
supporting assertion is informal and is not formalized; see
Appendix~\ref{app:sources}. Define $r=0$ on the null event of
singularity. Then $r$ is Borel, and almost surely
\begin{equation}\label{eq:reserve-identities}
 \mathbf1^{\mathsf T}r=0,\qquad H^{\mathsf T}r=b.
\end{equation}
Moreover,
\begin{equation}\label{eq:reserve-bounds}
 \norm r_2=\OP(d^{-1/2}),\qquad
 \left\|\sum_{j=1}^k r_jH_jH_j^{\mathsf T}\right\|_{\mathrm{op}}
       =\OP\big(\sqrt{\log(2d)}\big)=o_{\Pbb}(\sqrt d).
\end{equation}
\end{lemma}

\begin{proof}
The identities in~\eqref{eq:reserve-identities} follow directly from the
definition. The matrix $K$ has the law of the Gram matrix of $k-1$
independent Gaussian rows. Since $k-1\ge d$ eventually, it is invertible
almost surely. This step is informal, and the formal proof does not use
it. That proof shows instead that the event on which $K$ is singular or
the correction is large has probability tending to zero: this is
\path{reserve_correction_converges} in
\path{library/random_matrix/reserve_convergence.v}. On the event
$\opnorm{k^{-1}K-I_d}\le1/2$, the lemma \path{centered_correction_bounds}
gives the invertibility of $K$ and~\eqref{eq:reserve-identities}. Write
\[
 \overline H=\frac1k\sum_jH_j,\qquad B_j=H_jH_j^{\mathsf T}-I_d,
 \qquad \mathcal B^*(t)=\sum_jt_jB_j.
\]
Gaussian covariance concentration and the Gaussian norm bounds give
\begin{align}
 \opnorm{k^{-1}K-I_d}&=\OP(d^{-1/2}),&
 \opnorm H&=\OP(d),\label{eq:reserve-covariance}\\
 \norm{\overline H}&=\OP(d^{-1/2}),&
 \opnorm{\textstyle\sum_jB_j}&=\OP(d^{3/2}).\nonumber
\end{align}
It follows that
\begin{equation}\label{eq:reserve-inverse}
 \norm v=\OP(d^{-3/2}),\qquad
 \left\|v-\frac bk\right\|=\OP(d^{-2}),\qquad
 \norm r_2^2=b^{\mathsf T}K^{-1}b=\OP(d^{-1}).
\end{equation}

We also use the quadratic-feature bound
\begin{equation}\label{eq:reserve-adjoint}
 \norm{\mathcal B^*(t)}_F\le Ld\,\norm t_2
       \qquad\text{for every }t\in\R^k,
\end{equation}
on events of probability tending to one, for a fixed $L$ depending on
$\gamma$. This is the Gaussian case of \path{mixture_adjoint_bound},
with the factor $\sqrt d$ restored because $B_j$ is unnormalized;
its analytic source is \citet[Proposition 3.1]{KoehlerSohn2026}.

Let $D=\sum_jr_jH_jH_j^{\mathsf T}$. Since $\sum_jr_j=0$, equations
\eqref{eq:reserve-covariance}--\eqref{eq:reserve-adjoint} give
\begin{align}
 D&=D_0+
   \mathcal B^*\!\left(H\left(v-\frac bk\right)\right)
   -(\overline H^{\mathsf T}v)\sum_jB_j,
       \label{eq:reserve-decomposition}\\
 D_0&=\frac1k\sum_j(H_j^{\mathsf T}b)B_j.\nonumber
\end{align}
The middle term has operator norm at most
$Ld\,\opnorm H\norm{v-b/k}=\OP(1)$. The last term has operator norm
$\OP(d^{-1/2})$. Thus it suffices to bound $D_0$.

Represent $H_j=\varepsilon_jU_j$, with independent standard Gaussian
$U_j$ and independent Rademacher signs, all independent of $b$. We may
use an event $E_d$ depending only on $(U_jU_j^{\mathsf T})_j$, with
$\Pbb(E_d)\to1$, on which~\eqref{eq:reserve-adjoint} and
$\max_j\norm{U_j}^2\le L_0d$ hold for fixed constants $L,L_0$.
On this event, duality in~\eqref{eq:reserve-adjoint}, applied to
$uu^{\mathsf T}$ for $\norm u=1$, gives
\begin{equation}\label{eq:reserve-fourth}
 \sum_j(U_j^{\mathsf T}u)^4
 \le2\sum_j\big((U_j^{\mathsf T}u)^2-1\big)^2+2k
 \le2L^2d^2+2k\le L_1k.
\end{equation}
Consequently, for all $b,u\in\R^d$ with $\norm u=1$,
\begin{align*}
 \sum_j(U_j^{\mathsf T}b)^2(U_j^{\mathsf T}u)^2
     &\le L_1k\norm b^2,\\
 \sum_j(U_j^{\mathsf T}b)^2&\le\sqrt{L_1}\,k\norm b^2,
\end{align*}
by Cauchy--Schwarz and homogeneity. With $B_j=U_jU_j^{\mathsf T}-I_d$,
\[
 B_j^2=(\norm{U_j}^2-2)U_jU_j^{\mathsf T}+I_d
       \preceq L_0d\,U_jU_j^{\mathsf T}+I_d.
\]
The preceding estimates therefore yield the matrix variance bound
\begin{equation}\label{eq:reserve-variance}
 \left\|\sum_j(U_j^{\mathsf T}b)^2B_j^2\right\|_{\mathrm{op}}
       \le L_2dk\norm b^2.
\end{equation}

Conditional on $b$ and $(U_j)_j$, the matrix
$D_0=k^{-1}\sum_j\varepsilon_j(U_j^{\mathsf T}b)B_j$ is a Rademacher
series. The expected-norm bound of
\citet[Theorem 4.1.1]{TroppMatrixConcentration2015} and
\eqref{eq:reserve-variance} give, on $E_d$,
\begin{equation}\label{eq:reserve-rademacher}
 \E_\varepsilon\opnorm{D_0}
 \le\frac{\sqrt{2\log(2d)}}{k}
     \left\|\sum_j(U_j^{\mathsf T}b)^2B_j^2\right\|_{\mathrm{op}}^{1/2}
 \le L_3\norm b\sqrt{\frac{d\log(2d)}{k}}.
\end{equation}
The factor $\sqrt{2\log(2d)}$ is the constant of Tropp's theorem, which
is cited and not formalized. The Rocq proof uses
\path{rademacher_mx_series_variance} in
\path{library/random_matrix/rademacher_series.v}, which
gives~\eqref{eq:reserve-rademacher} with $40\sqrt{\log(2d)}$ in place of
$\sqrt{2\log(2d)}$. Only the constants $L_3$ and $L_4$ change, and the
conclusion $\opnorm{D_0}=\OP(\sqrt{\log(2d)})$ is the same.

For clarity, no moment assumption on the random vector $b$ is needed
beyond tightness of $\norm b/\sqrt d$. For $M,T>0$, conditional Markov's
inequality gives, for all sufficiently large $d$,
\[
 \Pbb\{\opnorm{D_0}>T\sqrt{\log(2d)}\}
 \le\Pbb(E_d^c)+\Pbb\{\norm b>M\sqrt d\}+\frac{L_4M}{T}.
\]
First choose $M$ large and then $T$ large. This proves
$\opnorm{D_0}=\OP(\sqrt{\log(2d)})$. Equation~\eqref{eq:reserve-decomposition}
proves~\eqref{eq:reserve-bounds}. All constants above are independent of
$b$ and $d$, and may depend on $\gamma$.
\end{proof}

\subsection{Nonexistence above \texorpdfstring{$1/4$}{1/4}, uniformly over translations}

Suppose $\alpha>1/4$ and fix $\beta\in(1/4,\alpha)$. Split the sample
into its first $N=\lfloor\beta d^2\rfloor$ points and its remaining
$k=n(d)-N$ points. Thus $k/d^2\to\alpha-\beta>0$.
Apply Lemma~\ref{lem:psd-certificate} to the first part. It supplies
$y$, $M_0$, and $b=\sum_{i\le N}y_iZ_i$, with
$\norm b=\OP(\sqrt d)$ and
$M_0\succeq c_\beta\sqrt d\,I_d$ with probability tending to one.
The reserve matrix $H$, with rows $Z_{N+j}^{\mathsf T}$, is independent
of $b$. Lemma~\ref{lem:reserve-correction} supplies $r$ with
$H^{\mathsf T}r=b$, $\sum_jr_j=0$, and
$\opnorm D=o_{\Pbb}(\sqrt d)$, where
$D=\sum_jr_jZ_{N+j}Z_{N+j}^{\mathsf T}$.

Define weights on the whole sample by
$\lambda=(y_1,\ldots,y_N,-r_1,\ldots,-r_k)$. They satisfy
\begin{equation}\label{eq:translation-certificate}
 \sum_i\lambda_i=0,\qquad \sum_i\lambda_iZ_i=0,\qquad
 M:=\sum_i\lambda_iZ_iZ_i^{\mathsf T}=M_0-D
       \succeq\frac{c_\beta}{2}\sqrt d\,I_d
\end{equation}
with probability tending to one. The weights are Borel functions of the
sample, using the choices in the two lemmas. We may take $\Omega_d$ to
be the Borel event on which $K$ is invertible and the last inequality
in~\eqref{eq:translation-certificate} holds. The two equalities then
hold exactly on $\Omega_d$.

For every $m\in\R^d$, expansion of the outer products gives on this
same event
\begin{equation}\label{eq:translation-cancellation}
 \sum_i\lambda_i(Z_i+m)(Z_i+m)^{\mathsf T}=M.
\end{equation}
If a nonzero $Q\succeq0$ fitted the translated points at level one,
then
\[
 0=\sum_i\lambda_i
  =\sum_i\lambda_i(Z_i+m)^{\mathsf T}Q(Z_i+m)
  =\tr(QM)
  \ge\frac{c_\beta}{2}\sqrt d\,\tr Q>0,
\]
a contradiction. A fitting matrix is necessarily nonzero, so this
excludes all PSD fits. The certificate and the event do not depend on
$m$; the assertion is therefore simultaneous over every translation,
without a union bound or a bound on its norm. More generally, for every
nonzero $Q\succeq0$, every $\ell\in\R^d$, and every $c\in\R$,
\[
 \sum_i\lambda_i
   (Z_i^{\mathsf T}QZ_i+2\ell^{\mathsf T}Z_i+c)
 =\tr(QM)>0.
\]
Thus no quadric of the form~\eqref{eq:no-quadric} contains all the
points. This includes ellipsoids with arbitrary centers. Finally,
Remark~\ref{rem:sign-invariance} transfers both threshold conclusions
for centered fits to the sign mixtures and completes
Theorem~\ref{thm:unrestricted}.

\subsection{The fitting-error limit for every deterministic shift}
\label{sec:shifted-error}

The same certificate gives a quantitative residual bound uniformly
over translations. Retaining its full margin, and then letting the
reserve fraction decrease after taking the dimension limit, identifies
the limiting error exactly.

\begin{theorem}[The Gaussian fitting-error limit is unchanged]
\label{thm:shifted-error}
Let $n(d)/d^2\to\alpha>0$ and let $Z_i$ be independent standard Gaussian
points. For every deterministic sequence $m_d\in\R^d$, with no norm
restriction,
\begin{equation}\label{eq:shifted-error-limit}
 \Gamma_{n,d}(Z+m_d)\xrightarrow{\Pbb}e_\star(\alpha,3).
\end{equation}
Moreover, for every $\delta>0$,
\begin{equation}\label{eq:uniform-error-lower}
 \Pbb\left\{\inf_{m\in\R^d}\Gamma_{n,d}(Z+m)
                 \ge e_\star(\alpha,3)-\delta\right\}\longrightarrow1.
\end{equation}
The infimum over all translations also converges in probability to
$e_\star(\alpha,3)$. Equation~\eqref{eq:shifted-error-limit} holds for
the sign mixtures in Remark~\ref{rem:sign-invariance} as well.
\end{theorem}

\begin{proof}
First suppose $\alpha>1/4$ and fix $1/4<\beta<\alpha$. Use the split
and certificate of the preceding subsection. On events of probability
tending to one, $\Gamma_{N,d}(Z_1,\ldots,Z_N)>0$, so $\norm y_2=1$ and
\[
 \norm\lambda_2^2=1+\norm r_2^2=1+\OP(d^{-1}),\qquad
 \lambda_{\min}(M)\ge
    \sqrt{N/d}\sqrt{\Gamma_{N,d}(Z_1,\ldots,Z_N)}-\opnorm D.
\]
For every $m$ and every $R\succeq0$ with $\tr R=d$, the annihilation
identities give
\[
 \langle\lambda,\mathcal L_{Z+m}(R)\rangle
   =\frac{\tr(RM)}{\sqrt{nd}}
   \ge\sqrt{\frac dn}\lambda_{\min}(M).
\]
Cauchy--Schwarz and minimization over $R$ therefore give, simultaneously
for every $m$,
\begin{equation}\label{eq:sharp-uniform-error-bound}
 \sqrt{\Gamma_{n,d}(Z+m)}\ge
 \frac{\left[
    \sqrt{N/n}\sqrt{\Gamma_{N,d}(Z_1,\ldots,Z_N)}
       -\sqrt{d/n}\opnorm D\right]_+}{\norm\lambda_2}.
\end{equation}
For this fixed $\beta$, the square of the right side converges in
probability to $(\beta/\alpha)e_\star(\beta,3)$, because
$\opnorm D=\OP(\sqrt{\log(2d)})=\oP(\sqrt d)$.
Given $\delta>0$, continuity in~\eqref{eq:gaussian-error-limit} allows
us to choose a fixed $\beta<\alpha$ sufficiently close to $\alpha$ to
deduce~\eqref{eq:uniform-error-lower}. The choice is made before
letting $d\to\infty$; no estimate uniform in a vanishing reserve
fraction is needed. For $\alpha\le1/4$, the same conclusion follows
from $\Gamma\ge0$ and $e_\star(\alpha,3)=0$.

For the upper bound, put $p=d-1$ and choose a deterministic isometry
$J_d:\R^p\to\R^d$ with image orthogonal to $m_d$; if $m_d=0$, choose
any such isometry. The points $g_i=J_d^{\mathsf T}Z_i$ are independent
standard Gaussian vectors in $\R^p$. For every $A\succeq0$ with
$\tr A=p$, the matrix $R=(d/p)J_dAJ_d^{\mathsf T}$ is PSD of trace $d$,
and
\[
 (Z_i+m_d)^{\mathsf T}R(Z_i+m_d)=\frac dp\,g_i^{\mathsf T}Ag_i.
\]
Taking centered residuals with their respective normalizations yields
\begin{equation}\label{eq:projection-error-upper}
 \Gamma_{n,d}(Z+m_d)\le\frac dp\,\Gamma_{n,p}(g).
\end{equation}
Since $n/p^2\to\alpha$ and $d/p\to1$, the right side converges to
$e_\star(\alpha,3)$ by~\eqref{eq:gaussian-error-limit}.
Together with~\eqref{eq:uniform-error-lower}, this
proves~\eqref{eq:shifted-error-limit} at every $\alpha>0$, including
$\alpha=1/4$. It does not decide exact fitting at the critical density.
The infimum over translations is bounded above by $\Gamma_{n,d}(Z)$,
which has the same limit. It is measurable because $\Gamma$ is
continuous in the sample and the infimum may be taken over rational
translations. Finally, rowwise sign changes leave every quadratic
measurement unchanged, proving the assertion for sign mixtures.
\end{proof}

For a simpler positive lower bound at a fixed split, on $\Omega_d$,
\[
 \inf_{m\in\R^d}\Gamma_{n,d}(Z+m)
 \ge\frac{c_\beta^2d^2}{4n\norm\lambda_2^2}
 =\frac{c_\beta^2}{4\alpha}+\oP(1).
\]
Theorem~\ref{thm:shifted-error} strengthens this bound to the full
Gaussian limit.

\section{Symmetric mixtures beyond Gaussian data}
\label{sec:symmetric-mixtures}

We now specify a class for which the same certificate works. Let the
law of a point in $\R^d$ be a finite or countable mixture
\begin{equation}\label{eq:symmetric-mixture-law}
 \rho_d=\sum_s w_{d,s}\bigotimes_{j=1}^d\mu_{d,s,j},\qquad
 w_{d,s}\ge0,\quad\sum_s w_{d,s}=1.
\end{equation}
Assume that every component coordinate has mean zero, variance one,
and subgaussian norm at most a fixed $K_0$, and that, for a fixed
$\kappa>1$,
\begin{equation}\label{eq:mixture-fourth}
 \sum_s w_{d,s}\int t^4\,\mu_{d,s,j}(dt)=\kappa
       \qquad(1\le j\le d).
\end{equation}
These are the admissibility assumptions of the formal mixture theorem.
Assume in addition that $\rho_d$ is centrally symmetric: $Z$ and $-Z$
have the same law. This holds, in particular, if every
$\mu_{d,s,j}$ is symmetric. No independence of the coordinates under
the mixture itself is required.

Write $\alpha_\star(\kappa)$ for the unshifted mixture threshold. In
the notation of Section~1, it is
\[
 \alpha_\star(\kappa)
   =v(\omega_\kappa)+\frac{\kappa-3}{2}m(\omega_\kappa)^2,
 \qquad \omega_\kappa=\frac{\kappa-3}{2}m(\omega_\kappa),
\]
where the root is unique. The threshold lies in $(0,1/2)$ for every
$\kappa>1$.
Let $e_\star(\beta,\kappa)$ denote the unshifted squared-error limit
from the same mixture theorem; it is positive when
$\beta>\alpha_\star(\kappa)$.

\begin{theorem}[Translations of symmetric admissible mixtures]
\label{thm:symmetric-mixtures}
Let $Z_1,\ldots,Z_n$ be independent with law~\eqref{eq:symmetric-mixture-law}
under the preceding assumptions, and let $n/d^2\to\alpha>0$.
\begin{enumerate}
\item If $\alpha>\alpha_\star(\kappa)$, then on Borel events of
      probability tending to one there is no nonzero $Q\succeq0$,
      $\ell\in\R^d$, and $c\in\R$ with
      $Z_i^{\mathsf T}QZ_i+2\ell^{\mathsf T}Z_i+c=0$ for all $i$.
      In particular, no translated sample $Z+m$ has a PSD fit,
      simultaneously for all $m\in\R^d$. There is also a constant
      $\eta>0$, depending only on $\alpha,\kappa$, such that
      \[
       \Pbb\left\{\inf_{m\in\R^d}\Gamma_{n,d}(Z+m)\ge\eta\right\}
       \longrightarrow1.
      \]
\item If $\alpha<\alpha_\star(\kappa)$ and a deterministic shift $m_d$
      is supported on $h_d=o(d)$ coordinates, then $Z+m_d$ has a
      positive definite fit with probability tending to one.
      There is no restriction on $\norm{m_d}$.
\end{enumerate}
The centered-fitting conclusions also hold for the sign mixture
$w_d\,\operatorname{Law}(Z+m_d)+(1-w_d)\,\operatorname{Law}(Z-m_d)$, for every
$w_d\in[0,1]$.
\end{theorem}

\begin{proof}
\emph{Nonexistence.}
Fix $\alpha_\star(\kappa)<\beta<\alpha$, and split the sample into
$N=\lfloor\beta d^2\rfloor$ points and $k=n-N$ reserve points.
The mixture squared-error theorem gives
$\Gamma_{N,d}\to e_\star(\beta,\kappa)>0$ in probability. The compact
convex projection in Lemma~\ref{lem:psd-certificate} is deterministic
and applies without change. It gives a Borel vector $y$, determined
by the outer products, with $\sum_i y_i=0$, $\norm y_2\le1$, and
\[
 M_0\succeq c_\beta\sqrt d\,I_d
 \quad\text{with probability tending to one},\qquad
 0<c_\beta<\sqrt{\beta e_\star(\beta,\kappa)}.
\]

We verify the remaining probabilistic inputs. Conditional on a mixture
label, independence and the uniform scalar subgaussian bound give
\[
 \E\exp(tu^{\mathsf T}Z)\le
       \exp(C_{K_0}t^2\norm u^2),\qquad
 \E ZZ^{\mathsf T}=I_d.
\]
The bound is uniform in the label, so it also holds after mixing.
A fixed net of the unit sphere therefore gives
$\max_{i\le n}\norm{Z_i}^2\le C d$ with probability tending to one.
Central symmetry permits the representation $Z_i=\varepsilon_iU_i$,
where the $U_i$ have law $\rho_d$ and the signs are independent.
The conditional second-moment argument in
Lemma~\ref{lem:psd-certificate} gives $\norm b=\OP(\sqrt d)$ for
$b=\sum_{i\le N}y_iZ_i$.

For the reserve rows, the same linear subgaussian bound gives
\[
 \opnorm{k^{-1}H^{\mathsf T}H-I_d}
       =\OP\big(\sqrt{d/k}\big),\qquad
 \norm{\overline H}=\OP\big(\sqrt{d/k}\big).
\]
Indeed, for a fixed unit vector $u$, the independent variables
$(u^{\mathsf T}H_j)^2-1$ have uniformly bounded subexponential norms.
Bernstein's inequality followed by a union bound over a $1/4$-net
with at most $9^d$ points proves the covariance estimate. The mean
estimate follows from $\E\norm{\overline H}^2=d/k$.
Since
$K/k=H^{\mathsf T}H/k-\overline H\overline H^{\mathsf T}$,
all the rates in~\eqref{eq:reserve-covariance} follow. In particular,
$K$ is invertible with probability tending to one. Almost-sure
invertibility is not needed and is not asserted for discrete laws.

The finite-mixture adjoint bound \path{mixture_adjoint_bound} gives
\eqref{eq:reserve-adjoint} with a constant depending on $K_0$ and
$\alpha-\beta$. It also holds for countable mixtures: truncate the
component labels at tail mass at most $d^{-4}$ and replace the tail
by one admissible component. The finite law has the same uniform
subgaussian bound, and coupling the $k=O(d^2)$ labels changes a sample
with probability at most $kd^{-4}=o(1)$. The adjoint estimate does not
require a prescribed fourth moment, so this replacement is sufficient
for that estimate. The squared-error input for countable mixtures
comes directly from their formal theorem.

Define $v=K^{-1}b$ and $r=P_kHv$ on the invertibility event and set
$v=r=0$ otherwise. Then~\eqref{eq:reserve-identities} holds exactly
on that event, and~\eqref{eq:reserve-inverse} holds in probability.
The decomposition~\eqref{eq:reserve-decomposition}, the fourth-moment
estimate~\eqref{eq:reserve-fourth} obtained from the adjoint bound,
and the variance estimate~\eqref{eq:reserve-variance} are deterministic.
Central symmetry gives the same independent-sign representation of
the reserve rows, so~\eqref{eq:reserve-rademacher} and its conditional
Markov argument prove
\[
 \norm r_2=\OP(d^{-1/2}),\qquad
 \opnorm D=\OP\big(\sqrt{\log(2d)}\big)=\oP(\sqrt d).
\]
Thus $\lambda=(y,-r)$ has all the properties
in~\eqref{eq:translation-certificate}. The deterministic cancellation
and quadric argument prove nonexistence on a Borel event independent
of the translation. The bound following
Theorem~\ref{thm:shifted-error} gives any fixed
$0<\eta<c_\beta^2/(4\alpha)$, proving the residual assertion.

\emph{Existence for shifts on few coordinates.}
Write $h=h_d$, $p=d-h$, and, after a coordinate permutation,
$Z_i=(U_i,g_i)$ and $m_d=(\widetilde m_d,0)$, with
$a_d=\norm{\widetilde m_d}$. The case $h=0$ is the unshifted theorem.
The complementary points $g_i$ remain an admissible mixture in
dimension $p$, with the same $K_0,\kappa$, and $n/p^2\to\alpha$.
The unshifted mixture theorem supplies, with probability tending to
one, $A\succ0$ with $\tr A=p$, $\cond(A)\le D$ for a fixed $D$, and
$g_i^{\mathsf T}Ag_i=q>0$ independent of $i$.

Set $t_d=(1+a_d)^{-1}$ and
\[
 y_i=-\frac{t_d\big(\norm{U_i+\widetilde m_d}^2-h-a_d^2\big)}{\sqrt p}.
\]
Conditional component independence and variance one imply
$\operatorname{Cov}(U_{ij}^2,U_{i\ell}^2)=0$ for $j\ne\ell$.
Central symmetry makes the covariance between
$\norm{U_i}^2$ and $\widetilde m_d^{\mathsf T}U_i$ zero. Hence
\[
 \E y_i^2=\frac{t_d^2}{p}\big(h(\kappa-1)+4a_d^2\big),\qquad
 \frac{\norm y_2}{p}=\OP\left(\sqrt{\frac{h+1}{d}}\right)=\oP(1).
\]
Uniform conditional Bernstein bounds for the sums of squares, and
the linear subgaussian bound, also give
\[
 \norm y_\infty
 =\OP\left(\frac{\sqrt{h\log d}+\log d}{\sqrt d}\right),\qquad
 \frac{(\log p)^4}{\sqrt p}\norm y_\infty=\oP(1).
\]
Since $\alpha<\alpha_\star(\kappa)<1/2$, mixture interpolation applies
to $(g_i)$ and gives $\opnorm\Delta=\oP(1)$. It holds simultaneously
for every target, so no independence between $(U_i)$ and $(g_i)$ is
required. With probability tending to one,
$Q=\operatorname{diag}(t_dI_h,A+\Delta)$ is positive definite and has
the common quadratic level
\[
 q+t_d(h+a_d^2)+\tr\Delta+\sqrt p\,b_0.
\]
This level is positive because $Q\succ0$ and each $g_i\ne0$ on the
fitting event. Rescaling proves the existence assertion.
Finally, central symmetry permits exactly the sign coupling of
Remark~\ref{rem:sign-invariance}, proving the centered-fitting
assertions for every mixture weight.
\end{proof}

For example, mixing the Rademacher product law on $\{\pm1\}^d$ with
the standard Gaussian product law, with fixed Gaussian weight $0<\theta\le1$,
gives $\kappa=1+2\theta$. The theorem excludes every translated PSD
fit above $\alpha_\star(1+2\theta)$ and supplies positive definite fits
below it for shifts supported on $o(d)$ coordinates. The case
$\kappa=1$ is excluded: unshifted Rademacher data lie exactly on a
sphere at every sample size. Existence for arbitrary dense shifts of
these non-Gaussian mixtures is not established by the present argument.

\appendix
\section{Formal results and their relation to the proof}
\label{app:sources}
\begingroup\raggedright\setlength{\parindent}{17pt}

The conclusions of this note are seven public propositions of the Rocq
development, listed below with the seven theorems that prove them. The
formalization is released under the annotated tag \texttt{shifted} of
the repository, placed on 24 October 2026 on commit
\texttt{305747af48d4f75e26378888a8111d591145fa52}, the final commit of
step 18, the step that merged the extension onto phase 14 under the
repaired checker. The statement files were accepted at commit
\texttt{51c2e14d5c2d75469f8147901bc8f664d33d4b6c} on 16 October 2026
and are frozen: rule 36 of \path{scripts/check_sources.py} compares
their code with that commit. The project definitions that they unfold
to are either in frozen files or already protected by earlier checks,
or among the 29 that \path{check/ShiftedDefinitions.v} fixes by
characterizations proved by computation. Each theorem has an axiom
check file, such as \path{check/ShiftedConditionAssumptions.v}, that
compiles only if the theorem proves its frozen statement, and
\path{check/ShiftedIdentities.v} checks the seven together (rule 37);
snapshots print the frozen statements and the characterizations
against builds of the commits at which they were written, and identity
drivers fix the kernel names that they use (rule 38); the check files
and the machinery of the extension are pinned, byte for byte, to the
integration commit \texttt{cc6fc9d49c73f74e0bccdb7f7a7d5716ce02958a}
(rule 39). The modules compiled with Rocq 9.1.1 and passed axiom
audits, independent source reviews and recursive kernel checking. The
references below point to the modules at the tag \texttt{shifted}.

The distinction is between conclusions and proof routes: some supporting
claims in the text are stronger than the estimates used in Rocq, and
only the weaker estimates are formalized. The text marks two of these
claims as not formalized: the almost-sure invertibility of $K$ in
Lemma~\ref{lem:reserve-correction} and the constant $\sqrt{2\log(2d)}$
in~\eqref{eq:reserve-rademacher}. The formal proof needs neither. The
paragraphs below specify these two differences and several others. They
do not list every estimate of the text that the formal proof replaces
by a weaker one.

\paragraph{Public statements.}
The statement files are at the top of \path{extension/}, and the
theorem files are in \path{extension/Shifted/}.
\begin{itemize}
\item Theorem~\ref{thm:shifted}: \path{shifted_condition_statement} in
  \path{extension/ShiftedConditionStatement.v}, proved by
  \path{shifted_condition_theorem} in
  \path{extension/Shifted/ConditionTheorem.v}.
\item Theorem~\ref{thm:unrestricted}, part~1:
  \path{shifted_gaussian_statement} in
  \path{extension/ShiftedGaussianStatement.v}, proved by
  \path{shifted_gaussian_theorem} in
  \path{extension/Shifted/GaussianTheorem.v}.
\item Theorem~\ref{thm:unrestricted}, part~2:
  \path{shifted_quadric_statement} in
  \path{extension/ShiftedQuadricStatement.v}, proved by
  \path{shifted_quadric_theorem} in
  \path{extension/Shifted/QuadricTheorem.v}.
\item Theorem~\ref{thm:shifted-error}: \path{shifted_error_statement} in
  \path{extension/ShiftedErrorStatement.v}, proved by
  \path{shifted_error_theorem} in \path{extension/Shifted/ErrorTheorem.v}.
\item Theorem~\ref{thm:symmetric-mixtures}, part~1:
  \path{shifted_mixture_unsat_statement} in
  \path{extension/ShiftedMixtureUnsatStatement.v}, proved by
  \path{shifted_mixture_unsat_theorem} in
  \path{extension/Shifted/MixtureUnsatTheorem.v}.
\item Theorem~\ref{thm:symmetric-mixtures}, part~2:
  \path{shifted_mixture_sat_statement} in
  \path{extension/ShiftedMixtureSatStatement.v}, proved by
  \path{shifted_mixture_sat_theorem} in
  \path{extension/Shifted/MixtureSatTheorem.v}.
  This statement needs no symmetry.
\item Remark~\ref{rem:sign-invariance} and the sign-mixture sentences of
  Theorems~\ref{thm:unrestricted}, \ref{thm:shifted-error}
  and~\ref{thm:symmetric-mixtures}: \path{shifted_sign_statement} in
  \path{extension/ShiftedSignStatement.v}, proved by
  \path{shifted_sign_theorem} in \path{extension/Shifted/SignTheorem.v}.
\end{itemize}
Several statements give, for every $\epsilon>0$ and every sufficiently
large dimension, a Borel exceptional set of probability at most
$\epsilon$ outside which the conclusion holds. In this way they express
events that serve all translations, all quadrics or several bounds at
once. A diagonal choice gives the events $\Omega_d$ of the text, as in
Section~2.

\paragraph{Gaussian conclusions.}
The paths in this paragraph are under \path{extension/Shifted/}.
Theorem~\ref{thm:shifted} is \path{shifted_condition_statement}, and
\path{shifted_condition_theorem} proves it from
\path{gaussian_shift_condition_threshold} and
\path{gaussian_shift_optimal_condition} in \path{GaussianCondition.v}.
The norm hypothesis is exactly $\norm{m_d}/\sqrt d\to0$; no logarithmic
restriction remains. The optimum is extended-real valued, and its
infimum is not assumed attained. The measurability stated in Section~1
is part of \path{shifted_condition_statement} and
\path{shifted_gaussian_statement}. In \path{GaussianCondition.v},
\path{measurable_gaussian_shift_fit} proves it for the events
$\{Z+m\in\Fit_{n,d}(C)\}$ and
\path{measurable_gaussian_shift_condition_number} for $\Copt(Z+m)$, for
every $m\in\R^d$; the case $m=0$ gives the fitting events and $\Copt$
themselves.

For Theorem~\ref{thm:unrestricted}, part~1 is
\path{shifted_gaussian_statement} and part~2 is
\path{shifted_quadric_statement}.
\path{gaussian_arbitrary_shift_condition_bounds} in
\path{GaussianBounds.v} chooses both constants before the sample-size
and shift sequences. The upper event may depend on that deterministic
shift sequence. The lower bound in \path{ConditionLower.v} holds on one
event for all translations. Above $1/4$,
\path{gaussian_quadric_good_events} in \path{GaussianErrorLower.v}
proves \path{shifted_quadric_statement}. It chooses the exceptional
event before all $Q,v,c$, with $Q\succeq0$ and $Q\ne0$. Its linear term
is $v^{\mathsf T}x$; take $v=2\ell$ to recover the notation of the text.

Theorem~\ref{thm:shifted-error} is \path{shifted_error_statement},
proved from \path{gaussian_shift_error_limit} and
\path{gaussian_translation_error_limit} in \path{GaussianError.v}.
These include $\alpha=1/4$ and arbitrary deterministic shifts.
The latter result concerns the infimum over all translations.
For positive sample size and dimension, its measurability is proved in
\path{TranslationError.v}, using open strict sublevel sets of an
infimum of continuous functions. The all-dimension extension is
\path{measurable_translation_fitting_error_all} in
\path{GaussianErrorLower.v}. No minimizing translation is selected. The uniform lower tail is
\path{gaussian_translation_error_lower_tail} in
\path{GaussianErrorLower.v}. A zero error limit at $1/4$ is not a proof
of exact fitting at that density.

\paragraph{Mixtures and signs.}
Part~1 of Theorem~\ref{thm:symmetric-mixtures} is
\path{shifted_mixture_unsat_statement}, proved from
\path{finite_symmetric_quadric} and \path{countable_symmetric_quadric}
in \path{SymmetricQuadric.v}. Part~2 is
\path{shifted_mixture_sat_statement}, proved from
\path{finite_sparse_shift_theorem} and \path{countable_sparse_shift_theorem}
in \path{SparseShift.v}. The nonexistence results require symmetry of
the actual mixture point law, not of each component. Their positive
residual margin depends only on $\alpha,\kappa$ and is chosen before
the common subgaussian bound and the families.

The sparse-shift existence results are stronger than the theorem stated
in the text: they need no symmetry. For the normalized target used in
Section~\ref{sec:symmetric-mixtures}, \path{SparseTargets.v} proves
\[
 \E y_i^2\le \frac{8K^4h+8}{p}.
\]
It bounds the squared sum of the centered quadratic and linear parts
rather than cancelling their covariance. \path{SparseTargetBounds.v}
uses the sufficient maximum bound
$\norm y_\infty=\OP((\sqrt h+1)\log(p+1)/\sqrt p)$.
Consequently the interpolation contribution is of order
$(\sqrt h+1)(\log p)^5/p$, which tends to zero under $h/p\to0$.
There is no concealed condition such as $h(\log p)^a=o(p)$.

The complementary dimension $p(d)=d-h(d)$ need not be injective or
monotone. \path{SparseInputs.v} uses the varying-dimension transfer in
\path{library/analysis/family_reindexing.v}, allowing different laws
at repeated values of $p$. The construction in \path{SparseExistence.v}
uses separate head and complement marginals and a union bound. It never
assumes independence of these two blocks after mixing. The final
statements have no auxiliary baseline-family or target-tail hypothesis.

The public statement \path{shifted_sign_statement} collects the sign
results. Its first conjunct, the equality of the joint laws of the outer
products, holds over any real type. The finite-sample sign identities
are in \path{SignMixture.v} and
\path{library/probability/symmetric_shift.v}. They cover weights in the
closed interval $[0,1]$, including dimension-dependent weights in the
limit theorems. \path{MixtureSign.v} supplies the finite and countable
mixture consequences. They transfer centered fit probabilities and the
laws of the optimum and fitting error. They do not transfer the quadric
event with arbitrary linear terms, or an infimum over further
translations: these are not invariant under separate signs on the data.

\paragraph{Inputs and normalization.}
\path{GaussianInputs.v} proves standard-normal admissibility with $K=2$
and fourth moment three. \path{GaussianExports.v} specializes the
existing condition, interpolation, adjoint and error results.
\path{GaussianRows.v} connects flat samples, row tuples, matrices,
coordinate splits and rotations. These are proved law identities, not
identifications made by notation.

The extracted gap is \path{gaussian_condition_residual_gap} in
\path{ConditionGap.v}. Its proof retains the loss in
\path{af_value_residual} and bounds the empirical intercept without an
exact-fit hypothesis. The existing frozen condition theorem is unchanged.
The curve properties remain those of \path{extension/Condition/Curve.v};
\path{extension/ConditionKappa/Compatibility.v} identifies this curve
and its inverse with the fourth-moment-three specialization.

The identity \path{sample_residual_gamma} in
\path{ProjectionCertificate.v} identifies the squared closest-residual
norm with the frozen \path{universality_fitting_error}. In particular,
the centered squared loss has factor $1/(nd)$, the trace is $d$, and the
best intercept is the empirical mean of nonnegative quadratic
measurements. The restriction to nonnegative intercepts therefore
causes no change. \path{library/convex/image_projection.v} proves
continuity of the unique closest image point even when its matrix
preimage is not unique.

\paragraph{Reserve correction and proof departures.}
The deterministic correction and cancellation identities are in
\path{library/matrix/centered_correction.v},
\path{library/matrix/centered_quadratic_correction.v} and
\path{library/matrix/quadric_certificate.v}.
The reserve estimates are in
\path{library/random_matrix/symmetric_reserve.v} and
\path{library/random_matrix/reserve_convergence.v}; \path{ReserveApplications.v} discharges
their hypotheses for Gaussian and admissible mixture rows.
The countable adjoint estimate is proved separately by label truncation
in \path{MixtureAdjoint.v}. Row subgaussianity is not used as a substitute
for this quadratic-feature estimate.

The matrix Rademacher estimate used formally is
\path{rademacher_mx_series_variance} in
\path{library/random_matrix/rademacher_series.v}. For $d\times d$
matrices its factor is $40\sqrt{\log(2d)}$ instead of
$\sqrt{2\log(2d)}$. The file derives it from the existing Gaussian
matrix-series estimate and a comparison between signs and Gaussians,
rather than formalizing Tropp's matrix Laplace argument. The constant
$\sqrt{2\log(2d)}$
in~\eqref{eq:reserve-rademacher} is therefore informal. The formal
constant suffices for the reserve error to be $o_{\Pbb}(\sqrt d)$.
The condition lower bound uses a centered sample-covariance deviation
estimate for independent isotropic subgaussian rows, together with the
raw quadratic adjoint bound. No separate centered Wishart limit is proved.

Lemma~\ref{lem:reserve-correction} states that the Gaussian reserve Gram
matrix is invertible almost surely for large dimensions. This stronger
supporting assertion is marked informal in the text and has not been
formalized. Rocq defines the correction to be zero when the matrix is
singular, and \path{reserve_correction_converges} proves that the event
on which the matrix is singular or the correction is large has
probability tending to zero. That is enough for every theorem
of this note. Likewise, the main fitting-error argument fixes a positive
reserve density before taking the dimension limit; it needs no estimate
uniform in a reserve fraction tending to zero.

Two estimates of Section~\ref{sec:conditioned-existence} are also
replaced by weaker ones. The text uses $\max_i|\xi_i|=\OP(\sqrt{\log d})$
for~\eqref{eq:sat-linf}. The formal proof uses the event
$\max_i|\xi_i|\le\log p$, whose complement has probability at most
$2n/p^4$ once $\log p\ge8$ (\path{gaussian_coordinate_log_tail} in
\path{GaussianTargetBounds.v}); this still gives $S_d=\oP(1)$ when
$a_d=o(\sqrt d)$. The text also says that the common level is positive
because the observations are almost surely nonzero. The formal proof
uses instead that every point is nonzero on the Gaussian fitting event
(\path{gaussian_conditioned_point_nonzero} and
\path{gaussian_axis_point_nonzero} in \path{GaussianExistence.v}), so it
needs no almost-sure argument.

\paragraph{Verification scope.}
Each of the seven public theorems uses exactly the usual seven
classical axioms of the ellipsoid development: the three \path{boolp}
axioms \path{propositional_extensionality},
\path{functional_extensionality_dep} and
\path{constructive_indefinite_description};
\path{FunctionalExtensionality.functional_extensionality_dep};
\path{Epsilon.epsilon_statement}; and \path{sig_forall_dec} and
\path{sig_not_dec} of \path{ClassicalDedekindReals}.
\path{Classical_Prop.classic} is not among them. The new general library
lemmas use at most the three \path{boolp} axioms. This is not an
axiom-free formalization.
Source review checks the mathematical meanings and quantifiers;
compilation and kernel checking verify the proofs of those propositions.
The seven statements are frozen with the repaired checker (rule 36 of
\path{scripts/check_sources.py}; rules 37 to 39 cover the axiom
checks, the snapshots and the pin), and \texttt{make check} requires
their seven axiom reports, 37 reports in all; at the tag
\texttt{shifted} it passes in two fresh clones of commit
\texttt{305747af48d4f75e26378888a8111d591145fa52}. Dense non-Gaussian
shift existence and exact fitting at the critical density remain outside
these conclusions.
\par\endgroup

\bibliographystyle{plainnat}
\bibliography{statements}
\end{document}
